\chapter{Gestión de producto y desarrollo de nuevos productos}\label{ch:product-management}

% Alcance: ajuste producto-mercado, descubrimiento y priorización (RICE/stage-gate), hoja de ruta y
%   construir-comprar-asociar, lanzamiento y adopción de nuevos productos (Bass). Capítulo nuevo
%   (cierra la brecha PM/NPD). La empresa de referencia ES un producto; «qué construir» es anterior
%   a «qué cobrar» (\cref{ch:pricing}).
% Borrador según estilo editorial: prosa fluida + recuadros de ejemplo/admonición; ejemplos
%   trabajados con la SaaS de referencia.

\section{Ajuste producto-mercado}\label{sec:product-market-fit}
% complexity: 2

¿Qué significa que un producto encaje con su mercado? La pregunta precede a toda
conversación sobre crecimiento en este libro. Una empresa que escala la adquisición
antes de tener un genuino ajuste producto-mercado (\emph{PMF}) está quemando capital
para llenar un cubo con fugas: los clientes llegan, no se quedan, y la curva de
\emph{retention} (\cref{eq:retention}) cae abruptamente desde la primera cohorte.
El ajuste producto-mercado es la condición en la que un producto satisface una necesidad
real y recurrente de manera tan completa que los clientes se sentirían genuinamente
afectados si perdieran el acceso a él. Antes de que esa condición se cumpla, incrementar
el gasto en \emph{marketing} acelera el fracaso; una vez que se cumple, el boca a boca
orgánico comienza a contribuir al crecimiento y la economía de la adquisición pagada
mejora de forma sustancial.

El diagnóstico canónico fue articulado por Sean Ellis tras estudiar casi un centenar de
empresas tecnológicas en etapa temprana. Propuso una sola pregunta de encuesta con
formulación negativa: «¿Cómo se sentiría si ya no pudiera usar este producto?» La
formulación negativa elimina el sesgo de cortesía incrustado en los puntajes de
satisfacción. El umbral es directo: si menos del \qty{40}{\percent} de los usuarios
activos recientes responden «muy decepcionado», el producto aún no ha logrado el ajuste
con ningún segmento que pueda definir. Con \qty{40}{\percent} o más, la observación
empírica respalda el escalamiento.
\textcite{ellis2017hacking} documenta el trabajo de estudio comparativo detrás de este
umbral; \textcite{cagan2018inspired} lo integra en un \emph{framework} más amplio de
descubrimiento de producto.

El puntaje de la encuesta es un indicador adelantado, no una definición. El verdadero
indicador empírico del ajuste producto-mercado es la forma de la curva de retención de
cohortes de \cref{eq:retention}. Un producto con ajuste genuino muestra una curva de
supervivencia asintótica: la curva se aplana hacia una fracción estable y distinta de
cero en lugar de continuar declinando hacia cero. Una curva que nunca se aplana ---sin
importar cuántos usuarios reporten sentirse muy decepcionados--- significa que el producto
tiene un problema de valor central que ninguna inversión en crecimiento resolverá. Los
dos diagnósticos son complementarios: la encuesta detecta problemas con anticipación,
mientras que la curva de retención los confirma o refuta con evidencia conductual a lo
largo del tiempo. \textcite{ries2011} articula la lógica de construir-medir-aprender que
rige la velocidad con que un equipo puede iterar hacia el ajuste.

\begin{example}[Diagnóstico de PMF para la SaaS de referencia]
La retención a \qty{12}{\text{-month}} de la empresa de referencia, de \qty{92.5}{\percent}
(\(c = \qty{0.65}{\percent}/\text{mo}\)), está muy por encima del umbral empírico de PMF:
una tasa de supervivencia anual de \qty{92.5}{\percent} produce una curva asintótica, lo
que confirma que el producto tiene un ajuste genuino con su base de clientes actual. La
preocupación está aguas arriba: la tasa de conversión de prueba a pago es solo del
\qty{31}{\percent} (800 pruebas, 200 de pago, según el \emph{funnel}). Esta brecha no es
un fallo de PMF ---los clientes activados y de pago se quedan--- sino un fallo de
\emph{onboarding}. El valor no emerge con suficiente rapidez durante la prueba. La
implicación diagnóstica: instrumentar el «momento ajá» (en este producto, ejecutar un
primer informe), medir el tiempo hasta el primer informe para los usuarios en prueba, y
probar intervenciones que lo reduzcan. La alta retención entre los usuarios de pago
confirma que el ajuste es real; la baja conversión de prueba confirma que el camino
hacia ese ajuste está oculto.
\end{example}

\begin{admonition}{Advertencia}
El ajuste producto-mercado es específico del segmento, no del producto en su totalidad.
Un puntaje del \qty{40}{\percent} promediado en toda la base de usuarios puede enmascarar
un \qty{65}{\percent} para un segmento y un \qty{15}{\percent} para otro. Escalar hacia
el segmento equivocado porque el puntaje combinado parecía aceptable es la manera en que
las empresas construyen bases de clientes grandes y con alto \emph{churn} que parecen
sanas hasta el primer ciclo de renovación. Segmente la encuesta por canal de adquisición,
tamaño de empresa y caso de uso antes de extraer conclusiones.
\textcite{torres2021continuous} trata el descubrimiento específico por segmento como la
base de la práctica de descubrimiento continuo.
\end{admonition}

\section{Descubrimiento y priorización}\label{sec:product-discovery}
% complexity: 4

Una vez que un producto ha demostrado ajuste con un segmento definido, la pregunta pasa
de \emph{¿este producto crea valor?} a \emph{¿qué mejoras crean más valor, más
rápidamente?} Este es el problema de priorización, y está persistentemente distorsionado
por sesgos cognitivos: la voz más alta del cliente frente a la necesidad del usuario
mediano, el entusiasmo del ingeniero frente al impacto de mercado, la complejidad de
implementación frente a la simplicidad del beneficio. La primera disciplina de la gestión
de producto consiste en reemplazar la intuición sobre la importancia con una estimación
estructurada del valor esperado por unidad de esfuerzo.

El \emph{framework} RICE operacionaliza esa estimación. Cada iniciativa propuesta se
evalúa en cuatro dimensiones, combinadas en un único número comparable:

\begin{equation}\label{eq:rice}
  \mathrm{RICE} \;=\; \frac{R \;\times\; I \;\times\; C}{E},
\end{equation}
donde \(R\) es \emph{Reach} (el número de usuarios o transacciones afectados por
periodo), \(I\) es \emph{Impact} (un multiplicador escalonado: 3 para masivo, 2 para
alto, 1 para medio, 0.5 para bajo, 0.25 para mínimo), \(C\) es \emph{Confidence}
(certeza análoga a una probabilidad sobre las estimaciones de Reach e Impact: 1.0, 0.8
o 0.5), y \(E\) es \emph{Effort} en personas-mes. El numerador premia las iniciativas
que alcanzan a muchos usuarios con alta confianza de impacto material; el denominador
penaliza los compromisos de ingeniería de gran envergadura. La fortaleza de la fórmula
radica en obligar a estimaciones explícitas de los cuatro factores, lo que hace visibles
los supuestos detrás de cualquier clasificación basada en intuición.

La estructura de gobernanza más amplia para el desarrollo de nuevos productos de mayor
complejidad o intensivos en hardware es el proceso \emph{stage-gate} (\cref{fig:stage-gate}).
El \emph{framework} segmenta el recorrido de desarrollo en fases discretas (Descubrimiento,
Definición del alcance, Caso de negocio, Desarrollo, Pruebas, Lanzamiento), separadas por
puertas de decisión donde el liderazgo interfuncional emite una instrucción explícita de
Avanzar, Cancelar, Pausar o Reciclar. Cada puerta impide que la inversión se acumule en
una iniciativa que ha fallado una condición necesaria ---atractivo de mercado en la
puerta 2, viabilidad financiera en la puerta 3, factibilidad técnica en la puerta 4.
\textcite{cooper2017winning} documenta la investigación empírica detrás del proceso y sus
adaptaciones para organizaciones de software; \textcite{cagan2018inspired} describe cómo
los equipos de descubrimiento y entrega operan en carriles paralelos en lugar de una
secuencia estricta.

\begin{figure}[htbp]
  \centering
  \includegraphics[width=0.82\linewidth]{stage-gate}
  \caption{Proceso \emph{stage-gate}: cinco etapas separadas por puertas de decisión
    Avanzar/Cancelar/Pausar. Los rombos representan decisiones de control; los rectángulos
    representan etapas de trabajo interfuncional paralelo. Una decisión de «Cancelar» o
    «Pausar» en cualquier puerta impide compromisos adicionales de recursos hasta que se
    resuelva la condición bloqueante.}
  \label{fig:stage-gate}
\end{figure}

\begin{example}[Priorización RICE: SSO vs.\ exportación CSV]
La cartera de desarrollo contiene dos funcionalidades. La funcionalidad A es la integración
de inicio de sesión único (SSO), solicitada principalmente por prospectos empresariales;
la funcionalidad B es la exportación de datos en CSV, solicitada por los usuarios operativos
dentro de cada cuenta existente. Aplicando \cref{eq:rice} con estimaciones por trimestre:

\medskip
\noindent Funcionalidad A (SSO): \(R = 200\) prospectos, \(I = 3\) (alto ---elimina un
obstáculo de compras), \(C = 0.8\), \(E = 4\) personas-mes:
\[
  \mathrm{RICE}_A \;=\; \frac{200 \times 3 \times 0.8}{4} \;=\; 120.
\]
Funcionalidad B (exportación CSV): \(R = 500\) usuarios activos, \(I = 1\) (medio ---mejora
de conveniencia), \(C = 0.9\), \(E = 1\) persona-mes:
\[
  \mathrm{RICE}_B \;=\; \frac{500 \times 1 \times 0.9}{1} \;=\; 450.
\]
B puntúa \num{3.75}\(\times\) más alto. La implicación: el cuello de botella en el
\emph{funnel} de la empresa de referencia es la tasa de conversión de prueba a pago del
\qty{31}{\percent}, no el ciclo de ventas empresarial (véase \cref{sec:toc} para la
lógica de restricciones). La funcionalidad A añadiría compras empresariales sin fricción
---valioso a largo plazo--- pero no ataca la restricción. La funcionalidad B ofrece un
beneficio tangible de flujo de trabajo a cada usuario de pago actual y potencial,
alcanzando a \num{500} usuarios por una sola persona-mes. Hasta que no se corrija la
tasa de conversión de prueba a pago, construir funcionalidades empresariales es una
optimización prematura.
\end{example}

\begin{admonition}{Nota}
Los puntajes RICE no son comparables entre perfiles de esfuerzo cualitativamente distintos.
Una refactorización arquitectónica de cinco personas-mes y un cambio de interfaz de una
persona-mes comparten el mismo denominador, pero implican riesgos, reversibilidad y deuda
técnica incomparables. Use RICE para clasificar funcionalidades dentro de una banda de
esfuerzo aproximadamente similar; aplique la gobernanza del \emph{stage-gate} a las
iniciativas lo suficientemente grandes como para requerir una puerta formal de caso de
negocio. El \emph{framework} suprime sesgos; no reemplaza el criterio sobre lo que el
negocio necesita. \textcite{torres2021continuous} advierte que los \emph{frameworks}
aplicados mecánicamente pueden convertirse en un sustituto del pensamiento que
pretendían habilitar.
\end{admonition}

\section{Hoja de ruta y construir-comprar-asociar}\label{sec:roadmap}
% complexity: 2

Incluso después de que RICE ha clasificado la cartera, surge una pregunta de orden superior
para cada iniciativa importante: ¿debería la capacidad construirse internamente, comprarse
(mediante licencia o adquisición) o entregarse a través de una asociación? Esta decisión
de construir-comprar-asociar moldea la hoja de ruta del producto de manera más duradera
que cualquier elección individual de funcionalidad, porque determina si la empresa acumula
propiedad intelectual propia o compra tiempo al costo del control estratégico. El criterio
financiero es \cref{eq:npv}: la opción con el mayor valor presente neto de la contribución
de margen esperada, correctamente cargada por integración y costo continuo, es la preferida.
El criterio estratégico es la prueba VRIN de \cref{sec:rbv-vrin}: construir solo cuando la
capacidad resultante sea Valiosa, Rara, Inimitable y No sustituible, porque solo entonces
la inversión genera una ventaja competitiva duradera. Construir funcionalidad de uso común
internamente es destrucción de costos de oportunidad.

El \emph{framework} práctico mapea tres rutas contra cuatro dimensiones: control estratégico,
costo total de propiedad, tiempo de comercialización y dependencia del proveedor. Construir
internamente maximiza el control y la propiedad de la propiedad intelectual, pero implica
la entrega más lenta y el mayor riesgo de ejecución. Comprar o licenciar ofrece el tiempo
de comercialización más rápido al costo de quedar atado al proveedor y a la dependencia de
su hoja de ruta. Asociarse equilibra velocidad y capacidad con gobernanza compartida y
derechos de propiedad intelectual negociados. La elección correcta depende de la posición
de la capacidad en la taxonomía VRIN, no de la preferencia de ingeniería.

\begin{example}[Mercado de integraciones: construir vs.\ asociarse con Zapier]
La empresa de referencia evalúa cómo atender al \qty{35}{\percent} de prospectos que solicitan
integraciones con herramientas de terceros (sistemas de RR.HH., CRMs, gestores de proyectos).
Dos opciones: construir un mercado de integraciones nativo (costo de ingeniería \money{80000}
inicial, dos ingenieros durante cuatro meses, más mantenimiento continuo); o asociarse con
una plataforma de automatización de flujos de trabajo consolidada a costo inicial cero a
cambio de un reparto del \qty{30}{\percent} de los ingresos en cualquier actualización de
plan atribuida a las integraciones.

Aplicando \cref{eq:npv}: la construcción nativa requiere \(I_0 = \money{80000}\) hoy. Si las
integraciones reducen el \emph{churn} anual en \qty{1}{\percent} sobre la base ARR de
\money{4000000}, el beneficio anual de retención es \money{40000}. Capitalizado como una
perpetuidad creciente a la tasa de descuento del \qty{11}{\percent} de la empresa:
\(\money{40000}/(0.11 - 0.03) = \money{500000}\).
NPV de construir \(\approx \money{500000} - \money{80000} = \money{420000}\): claramente
positivo.

La ruta de asociación ahorra \money{80000} iniciales pero paga un reparto del \qty{30}{\percent}
de ingresos por actualizaciones atribuidas a perpetuidad. Más críticamente: la capa de
integración es lo que ancla a los clientes al modelo de datos y flujo de trabajo del
producto. Esa es precisamente la definición de un \emph{moat} por costo de cambio (véase
\cref{sec:network-effects}). Externalizar las integraciones a un tercero significa que la
fidelización se acumula en la plataforma del socio, no en la de la empresa. La lógica VRIN
lo confirma: una capa de integración propia es Valiosa (reduce el \emph{churn}), Rara
(pocos productos SaaS de mercado medio construyen en profundidad), Inimitable (el
conocimiento del flujo de trabajo se incorpora con el tiempo) y No sustituible. La opción
de construir gana ---no solo por NPV sino por razones estratégicas.
\end{example}

\begin{admonition}{Advertencia}
Construir funcionalidad de uso común ---características que cualquier equipo de ingeniería
podría replicar en un fin de semana, que no requieren datos propietarios y que no generan
ningún efecto de red--- es destrucción de costos de oportunidad. Cada persona-mes destinada
a replicar lo que puede licenciarse a bajo costo es una persona-mes que no se destina a la
capacidad que constituye la diferenciación real de la empresa. La prueba es siempre VRIN,
aplicada antes de redactar el plan del sprint.
\end{admonition}

\section{Lanzamiento y adopción de nuevos productos}\label{sec:adoption-curves}
% complexity: 4

Construir el producto correcto y fijarlo al precio adecuado no garantiza la adopción. La
velocidad a la que un mercado absorbe un nuevo producto está gobernada por dinámicas
sociales que las decisiones individuales de producto pueden influir pero no anular. El
Modelo de Difusión de Bass es la descripción matemática fundacional de cómo se desarrollan
esas dinámicas. Descompone la adopción en dos fuerzas conductuales: la \emph{innovación}
---individuos que adoptan espontáneamente, impulsados por medios externos y una afinidad
intrínseca por la nueva tecnología--- y la \emph{imitación} ---individuos que adoptan
porque observan a sus pares que ya lo han hecho. La razón entre estas dos fuerzas
determina no solo el nivel de adopción final, sino el momento y la altura del pico.

En la forma de tiempo discreto utilizada para la elaboración de pronósticos prácticos, los
nuevos adoptantes en el periodo \(t\) son:
\begin{equation}\label{eq:bass}
  N(t) \;=\; \bigl[p + q\,\tfrac{N(t-1)}{M}\bigr]\,\bigl[M - N(t-1)\bigr],
\end{equation}
donde \(M\) es el mercado total direccionable (techo máximo de adoptantes), \(p\) es el
coeficiente de innovación, \(q\) es el coeficiente de imitación y \(N(t-1)\) es la base
instalada acumulada al inicio del periodo \(t\). El primer término entre corchetes es el
riesgo instantáneo de adopción: la probabilidad de que un no adoptante convierta,
combinando una tasa base (\(p\)) con un término de presión social que crece conforme
aumenta la base instalada (\(q \cdot N(t-1)/M\)). El segundo término entre corchetes es
el conjunto de no adoptantes restantes. Cuando la base instalada es pequeña, el término
social es insignificante y el crecimiento es lento; a medida que la base se acumula, el
término social acelera la adopción hasta que la saturación del mercado comprime el conjunto
restante y el crecimiento desacelera. El resultado es la característica curva de adopción
acumulada en forma de S y la curva de nuevos adoptantes en forma de campana que se muestran
en \cref{fig:bass-diffusion}. Las estimaciones de metaanálisis en diversas categorías de
productos sitúan \(\bar{p} \approx 0.03\) y \(\bar{q} \approx 0.38\)
\parencite{bass1969newproduct}.

\begin{figure}[htbp]
  \centering
  \includegraphics[width=0.86\linewidth]{bass-diffusion}
  \caption{Modelo de difusión de Bass: curvas de adopción acumulada en forma de S para el
    lanzamiento del nivel empresarial (\(M = 500\), \(p = 0.02\), \(q = 0.3\)). La curva de
    \emph{solo innovadores} (coeficiente \(p\) actuando solo, sin imitación) sigue una
    difusión puramente lineal; la curva completa de Bass se curva en sigmoide una vez que
    la base instalada es lo suficientemente grande para que el boca a boca domine. La línea
    vertical discontinua marca el pico de nuevos adoptantes (mes 10,
    \(\approx\num{43}\) cuentas nuevas).}
  \label{fig:bass-diffusion}
\end{figure}

\begin{example}[Lanzamiento del nivel empresarial: calibración de Bass]
La empresa de referencia lanza un nivel de precios empresarial dirigido a \(M = 500\) cuentas
de mercado medio en su segmento direccionable. La prospección saliente histórica sugiere una
tasa de conversión base de \(p = 0.02\) (innovadores ---cuentas que evalúan y compran de
forma independiente). Las entrevistas de éxito del cliente sugieren que las referencias boca
a boca dentro de las redes de pares del sector están activas, estimadas en \(q = 0.3\).

Aplicando \cref{eq:bass} de forma iterativa:

\medskip
\textbf{Mes 1.}\ Sin base instalada todavía. Nuevos adoptantes \(= (0.02 + 0.3 \times 0/500)
\times (500 - 0) = 0.02 \times 500 = 10\). Acumulado: \num{10}.

\textbf{Mes 2.}\ El término social contribuye ahora. Nuevos adoptantes \(\approx (0.02 + 0.3
\times 10/500) \times 490 = 0.026 \times 490 \approx 13\). Acumulado: \num{23}.

La aceleración continúa hasta el mes 10, donde el conteo de nuevos adoptantes alcanza su pico
en aproximadamente \num{43} por mes (acumulado \(\approx\num{268}\) de \num{500}). Después
del pico, el conjunto direccionable restante se contrae más rápido de lo que crece el término
social, y los conteos de nuevos adoptantes regresan hacia el piso de solo innovadores.

La implicación gerencial: el plan de lanzamiento debe concentrar el alcance en los meses 1--3,
asegurando a los primeros innovadores cuya presencia siembra el término social. Cada referencia
de cliente temprana ---caso de estudio, referido de un par, testimonio en conferencia del sector---
eleva el \(q\) efectivo y adelanta el pico. Este es precisamente el papel del producto mínimo
viable en el ciclo de la \emph{startup} lean (\cref{sec:lean-startup}): un MVP colocado con
cinco a diez adoptantes tempranos calibra empíricamente tanto \(p\) como \(q\) antes de la
inversión completa en el \emph{funnel} de lanzamiento.
\end{example}

\begin{admonition}{Advertencia}
El modelo de Bass describe la difusión a nivel macro en una población de adoptantes
potenciales independientes ---no es un modelo de CLV por cohorte. \Cref{eq:bass} no dice
nada sobre cuánto vale cada cuenta individual después de la adopción, ni sobre las
dinámicas de \emph{retention} dentro de la base instalada una vez adquirida. Los dos
modelos son complementarios: Bass pronostica la oferta de cuentas nuevas a lo largo del
tiempo; \cref{eq:retention} y \cref{eq:clv} en \cref{ch:customer-economics} determinan
el valor extraído de cada cuenta después de que se incorpora. Confundir los dos ---usar
las curvas de adopción de Bass para proyectar ingresos sin un análisis separado de
retención por cohorte--- es el error de modelado más común en la planificación financiera
de la etapa de lanzamiento. \textcite{bass1969newproduct} es explícito en cuanto al
alcance: el modelo rige la difusión de la primera compra, no la compra repetida ni la
retención.
\end{admonition}
